\section{Conclusion}

This paper described a minimal framework for biological systems using information density \(\I\), information flow \(\J\), and induced potential \(\PhiI\). Within this framework, development is interpreted as convergence toward stable structure, life as sustained stable motion, and intelligence as a second attractor supported by life.

The main point is that biological form can be described without treating design or purpose as primitive explanatory terms. The paper does not aim to reproduce every biological mechanism in detail. It extracts a common structure across biological phenomena.

The framework uses a small set of elements: state, flow, induced bias, and stability. These elements allow development, morphogenesis, life maintenance, probability revision, and the later emergence of intelligence to be placed in one descriptive form.

Several questions remain open. The parameter set \(\theta\) must be specified for concrete systems. The empirical forms of \(\I\), \(\J\), and \(\PhiI\) must be identified. The implementation of intelligence requires additional modeling.

The result of the paper is therefore not intended as a final theory. It is a form of description. Its usefulness depends on future applications and refinements.
